% Mô tả chi tiết kiến trúc front_qconv4x3_proj10 (level32 + encoder quantum qconv4x3), MNIST 16x16.
% Viết từ mã nguồn full_quantum_16x16_level32_qconv (models/preprocessing.py, front.py, encoding.py,
% fusion_model.py, tn_core.py, reupload.py, unitary.py, pauli.py; run_level32_front.py).
% Biên dịch: pdflatex kien_truc_qconv4x3.tex (2 lần để có mục lục). Overleaf: compiler pdfLaTeX.
% Hình: nếu có hinh_kien_truc_qconv4x3.pdf (biên dịch từ hinh_kien_truc_qconv4x3.tex) cùng thư mục thì được chèn tự động.
\documentclass[11pt,a4paper]{article}
\usepackage[T5]{fontenc}
\usepackage[utf8]{inputenc}
\usepackage[vietnamese]{babel}
\usepackage[margin=2.2cm]{geometry}
\usepackage{amsmath,amssymb,bm}
\usepackage{booktabs,array}
\usepackage{graphicx}
\usepackage{enumitem}
\usepackage[hidelinks]{hyperref}

\newcommand{\ket}[1]{\lvert #1 \rangle}
\newcommand{\bra}[1]{\langle #1 \rvert}
\newcommand{\expval}[1]{\langle #1 \rangle}
\newcommand{\RR}{\mathbb{R}}
\newcommand{\CC}{\mathbb{C}}
\newcommand{\Tr}{\operatorname{Tr}}
\newcommand{\figfile}{hinh_kien_truc_qconv4x3.pdf}
\newcommand{\chenhinh}[2]{%
  \IfFileExists{\figfile}{%
    \begin{figure}[htbp]\centering
      \includegraphics[page=#1,width=\textwidth,height=0.42\textheight,keepaspectratio]{\figfile}
      \caption{#2}
    \end{figure}}{%
    \begin{center}\fbox{\parbox{0.9\textwidth}{\small (Hình: biên dịch \texttt{hinh\_kien\_truc\_qconv4x3.tex}
    rồi đặt \texttt{hinh\_kien\_truc\_qconv4x3.pdf} cạnh file này; trang #1: #2)}}\end{center}}}

\title{Kiến trúc \texttt{front\_qconv4x3\_proj10}:\\
QCNN cây tensor network với encoder quanvolution lượng tử,\\ MNIST $16\times16$}
\author{}
\date{15/11/2026}

\begin{document}
\maketitle

\begin{abstract}
Mô hình phân loại ảnh MNIST thu nhỏ về $16\times16$ gồm ba phần: (i) encoder lượng tử \emph{qconv4x3}:
ba tầng mạch 9 qubit trượt trên cửa sổ $3\times3$ của từng pixel, đo các kỳ vọng Pauli một qubit và cập nhật
dư hai góc mã hóa $(\theta,\phi)$ của pixel; (ii) ba cây QCNN tensor network, mỗi cây nhận 256 qubit (một qubit mỗi
pixel) theo một thứ tự quét khác nhau, nén về 8 qubit và đo thêm tầng 32 qubit; (iii) head cổ điển nhỏ.
Tổng $17.264$ tham số, trong đó $5.148$ ($29{,}8\%$) là tham số mạch. Mọi đại lượng lượng tử được tính chính xác
bằng mô phỏng; mô hình mô phỏng cổ điển hiệu quả, nên không có tuyên bố lợi thế lượng tử.
\end{abstract}

\tableofcontents

% =====================================================================================================
\section{Tổng quan}

\begin{equation*}
\underbrace{\text{ảnh } 28{\times}28}_{\text{MNIST}}
\xrightarrow{\text{aug, resize}} \bm{x}\in[0,1]^{256}
\xrightarrow[\text{3 tầng}]{\text{qconv4x3}} (\bm{\theta},\bm{\phi})\in\RR^{256\times2}
\xrightarrow{\text{mã hóa}} \bigotimes_{p=1}^{256}\ket{q_p}
\xrightarrow[\text{3 cây}]{\text{QCNN}} \bm{f}_t\in\RR^{372}
\xrightarrow{\text{head}} \text{10 logit}.
\end{equation*}

\begin{table}[htbp]\centering
\caption{Số tham số (khớp số runner in ra).}
\begin{tabular}{lrrl}
\toprule
Phần & Tham số & Loại & Ghi chú \\
\midrule
Encoder qconv4x3: góc mạch & 324 & lượng tử & $3$ tầng $\times$ $4$ lớp $\times$ $9$ qubit $\times$ $3$ góc \\
Encoder qconv4x3: readout & 168 & cổ điển & $3 \times$ Linear$(27,2)$ \\
3 cây QCNN & 4.824 & lượng tử & $3\times1.608$ \\
Head & 11.948 & cổ điển & $3\times3.750 + 698$ \\
\midrule
\textbf{Tổng} & \textbf{17.264} & & lượng tử $5.148$ ($29{,}8\%$) \\
\bottomrule
\end{tabular}
\end{table}

\chenhinh{1}{Luồng dữ liệu tổng quát. Xanh: mạch lượng tử (mô phỏng chính xác); cam: cổ điển.}

% =====================================================================================================
\section{Dữ liệu và tiền xử lý}
\label{sec:pre}

\subsection{Dữ liệu}
MNIST train $60.000$ ảnh, chia cố định (hoán vị với \texttt{split-seed} $20260911$) thành $50.000$ ảnh train và
$10.000$ ảnh validation. Bộ test chính thức $10.000$ ảnh chỉ dùng một lần cho mô hình chốt cuối.
Batch $64$.

\subsection{Chuyển ảnh thành tensor xám}
Mỗi ảnh MNIST là ảnh xám 8 bit $I\in\{0,\dots,255\}^{28\times28}$. \texttt{ToTensor} chia cho 255:
\begin{equation}
X = I/255 \in [0,1]^{1\times28\times28},
\end{equation}
một kênh, $0$ là nền đen, $1$ là nét trắng; không chuẩn hóa trung bình/độ lệch.
Hàm \texttt{rgb\_to\_grayscale} trong \texttt{preprocessing.py} chỉ có tác dụng với ảnh 3 kênh
($0{,}2989R+0{,}5870G+0{,}1140B$); với MNIST (1 kênh) nó giữ nguyên.

\subsection{Tăng cường dữ liệu (chỉ khi train)}
Trước \texttt{ToTensor}, trên ảnh $28\times28$: \texttt{RandomAffine} với góc xoay ngẫu nhiên đều trong
$[-5^\circ,5^\circ]$, dịch ngẫu nhiên tối đa $\pm0{,}05$ kích thước ảnh mỗi chiều ($\pm1{,}4$ pixel),
nội suy song tuyến, nền điền $0$. Ảnh validation và test không tăng cường.

\subsection{Thu nhỏ về $16\times16$}
Trong mô hình (\texttt{preprocess\_full\_quantum}), ảnh được nội suy song tuyến
(\texttt{F.interpolate}, \texttt{align\_corners=False}) về $16\times16$. Pixel đích $(i,j)$,
$i,j\in\{0,\dots,15\}$, lấy mẫu tại tọa độ nguồn
\begin{equation}
u_i = (i+\tfrac12)\cdot\tfrac{28}{16}-\tfrac12 = 1{,}75\,i+0{,}375, \qquad v_j = 1{,}75\,j+0{,}375,
\end{equation}
và
\begin{equation}
\tilde X_{ij} = \sum_{a\in\{\lfloor u_i\rfloor,\lfloor u_i\rfloor+1\}}\ \sum_{b\in\{\lfloor v_j\rfloor,\lfloor v_j\rfloor+1\}}
(1-|u_i-a|)(1-|v_j-b|)\,X_{ab}.
\end{equation}
Sau đó kẹp về $[0,1]$ và làm phẳng theo hàng:
\begin{equation}
x_p = \min(1,\max(0,\tilde X_{ij})), \qquad p = 16\,i + j \in\{0,\dots,255\}.
\end{equation}
Kết quả là vector $\bm{x}\in[0,1]^{256}$, mỗi phần tử là độ xám của một pixel.

% =====================================================================================================
\section{Encoder lượng tử qconv4x3}
\label{sec:enc}

\subsection{Trạng thái của encoder}
Mỗi pixel $p$ mang hai số thực $(\theta_p,\phi_p)$, khởi đầu
\begin{equation}
\theta^{(0)}_p = x_p, \qquad \phi^{(0)}_p = 0.
\end{equation}
Encoder gồm $S=3$ tầng nối tiếp; tầng $s$ biến $(\bm\theta^{(s-1)},\bm\phi^{(s-1)})$ thành
$(\bm\theta^{(s)},\bm\phi^{(s)})$. Mỗi tầng có tham số riêng, nhưng trong một tầng mọi pixel dùng chung tham số
(giống một lớp tích chập).

\chenhinh{2}{Một tầng encoder: cửa sổ $3\times3$, mạch 9 qubit 4 lớp, đo Pauli, readout và cập nhật dư.}

\subsection{Cửa sổ $3\times3$}
Với pixel $p$ ở hàng $i$, cột $j$, cửa sổ gồm 9 pixel, đánh số theo hàng:
\begin{equation}
\mathcal N(p) = \{(i+r,\, j+c) : r,c\in\{-1,0,1\}\}, \qquad k = 3(r+1)+(c+1)\in\{0,\dots,8\},
\end{equation}
pixel giữa là $k=4$. Pixel nằm ngoài ảnh nhận $\theta=\phi=0$ (đệm 0, tương ứng qubit $\ket0$).

\subsection{Mã hóa 9 qubit}
Pixel thứ $k$ của cửa sổ thành một qubit
\begin{equation}
\ket{e(\theta_k,\phi_k)} = R_Z(\pi\phi_k)R_Y(\pi\theta_k)\ket0
\simeq \cos\tfrac{\pi\theta_k}{2}\ket0 + e^{i\pi\phi_k}\sin\tfrac{\pi\theta_k}{2}\ket1
\label{eq:enc}
\end{equation}
($\simeq$: bằng nhau sai khác pha toàn cục). Đây là tham số hóa chuẩn của một qubit thuần trên mặt cầu Bloch,
với góc cực $\pi\theta$ và góc phương vị $\pi\phi$; vector Bloch là
$(\sin\pi\theta\cos\pi\phi,\ \sin\pi\theta\sin\pi\phi,\ \cos\pi\theta)$.
Pixel đen ($\theta=0$) là cực bắc $\ket0$, pixel trắng ($\theta=1$) là cực nam $\ket1$.
Trạng thái đầu vào của cửa sổ là tích
\begin{equation}
\ket{\psi_p^{\text{in}}} = \bigotimes_{k=0}^{8} \ket{e(\theta_k,\phi_k)} \in \CC^{2^9}.
\end{equation}

\subsection{Mạch biến phân 4 lớp}
Phép quay một qubit tổng quát (góc Euler $Z$-$Y$-$Z$):
\begin{equation}
R(a,b,c) = R_Z(a)\,R_Y(b)\,R_Z(c) =
\begin{pmatrix}
\cos\frac b2\, e^{-i(a+c)/2} & -\sin\frac b2\, e^{-i(a-c)/2}\\[2pt]
\sin\frac b2\, e^{i(a-c)/2} & \cos\frac b2\, e^{i(a+c)/2}
\end{pmatrix}.
\end{equation}
Lớp vướng víu là tích các cổng CZ trên 12 cạnh của lưới $3\times3$
(6 cạnh ngang, 6 cạnh dọc, không có cạnh chéo):
\begin{equation}
E = \{(0,1),(1,2),(3,4),(4,5),(6,7),(7,8),(0,3),(3,6),(1,4),(4,7),(2,5),(5,8)\},
\qquad
\mathrm{CZ}_E \ket{b_0\dots b_8} = (-1)^{\sum_{(j,k)\in E} b_j b_k}\ket{b_0\dots b_8}.
\end{equation}
Mạch của tầng $s$, với $L=4$ lớp và góc $\alpha^{(s)}_{\ell,k}=(a,b,c)$ riêng cho từng lớp $\ell$ và qubit $k$:
\begin{equation}
U^{(s)} = \prod_{\ell=4}^{1}\Big[\mathrm{CZ}_E \cdot \bigotimes_{k=0}^{8} R\big(\alpha^{(s)}_{\ell,k}\big)\Big],
\qquad
\ket{\psi_p^{\text{out}}} = U^{(s)}\ket{\psi_p^{\text{in}}}.
\end{equation}
(Tích viết từ phải sang trái: lớp 1 tác dụng trước.) Số góc mỗi tầng: $4\times9\times3=108$,
khởi tạo ngẫu nhiên đều trong $[-\pi,\pi]$.

\paragraph{Ý nghĩa của CZ (ví dụ tính bằng code).}
Khi mọi góc quay bằng 0 và chỉ có một lớp CZ, kỳ vọng $X$ của qubit giữa là
\begin{equation}
\expval{X_4} = \sin(\pi\theta_4)\cos(\pi\phi_4)\prod_{j\in\{1,3,5,7\}}\cos(\pi\theta_j),
\end{equation}
tức là qubit giữa ``nghe'' được bốn hàng xóm nối với nó: hàng xóm tối nhân $+1$, sáng nhân $-1$, xám $0{,}5$ nhân $0$.
Với $\theta_4=0{,}5$: hàng xóm đều tối cho $\expval{X_4}=+1$; chỉ hàng xóm trái sáng cho $-1$; hàng xóm trái
xám cho $0$ (tính bằng \texttt{QuantumConvStage.features}). Nét dọc và nét ngang (hai hàng xóm sáng) đều cho
$+1$; sau 4 lớp với góc khác 0, hai trường hợp này cho vector đặc trưng khác nhau (khoảng cách $1{,}46$ với góc
ngẫu nhiên lúc khởi tạo). Phép quay đổi trục, cho phép thông tin pha do CZ tạo ra hiện lên ở các phép đo,
và làm mạch phân biệt được hướng.

\subsection{Đặc trưng đo}
Từ trạng thái ra, tính chính xác kỳ vọng Pauli một qubit của cả 9 qubit:
\begin{equation}
\bm g_p = \big(\expval{X_k},\expval{Y_k},\expval{Z_k}\big)_{k=0}^{8} \in [-1,1]^{27},
\qquad \expval{O_k} = \bra{\psi_p^{\text{out}}} O_k \ket{\psi_p^{\text{out}}}.
\end{equation}
Trong code (\texttt{bloch\_of\_qubits}): với $\rho^{(k)}_{01}=\sum a_{\cdot0\cdot}\,\overline{a_{\cdot1\cdot}}$
(tổng trên các biên độ có qubit $k$ bằng 0 và bằng 1),
$\expval{X_k}=2\operatorname{Re}\rho^{(k)}_{01}$, $\expval{Y_k}=-2\operatorname{Im}\rho^{(k)}_{01}$,
$\expval{Z_k}=\sum|a_{\cdot0\cdot}|^2-\sum|a_{\cdot1\cdot}|^2$. Không đo tương quan nhiều qubit.
Trên phần cứng thật cần 3 cấu hình đo (cơ sở $Z$; $X$ qua Hadamard; $Y$ qua $S^\dagger$ rồi Hadamard),
mỗi cấu hình nhiều lần đo để ước lượng.

\subsection{Readout và cập nhật dư}
Một lớp tuyến tính dùng chung cho mọi pixel (như tích chập $1\times1$):
\begin{equation}
\begin{pmatrix}\Delta\theta_p\\ \Delta\phi_p\end{pmatrix} = W^{(s)}\bm g_p + \bm b^{(s)},
\qquad W^{(s)}\in\RR^{2\times27},\ \bm b^{(s)}\in\RR^2\ \ (56 \text{ tham số}),
\end{equation}
\begin{equation}
\theta^{(s)}_p = \theta^{(s-1)}_p + \Delta\theta_p, \qquad \phi^{(s)}_p = \phi^{(s-1)}_p + \Delta\phi_p.
\end{equation}
Lý do cần lớp này: mỗi pixel phải được mã hóa lại thành \emph{một} qubit thuần, tức đúng hai số $(\theta,\phi)$,
trong khi phép đo cho 27 số; lớp tuyến tính học tổ hợp nào của 27 số là có ích.
$W^{(s)},\bm b^{(s)}$ khởi tạo bằng 0, nên lúc bắt đầu train encoder là ánh xạ đồng nhất và toàn mô hình trùng với
level32 không encoder. Lớp này không trộn các pixel với nhau; việc trộn hàng xóm do mạch đảm nhiệm.

\subsection{Ba tầng và vùng nhìn}
Tầng $s$ dùng $(\bm\theta^{(s-1)},\bm\phi^{(s-1)})$ đã cập nhật làm đầu vào. Giữa các tầng là đo rồi mã hóa lại:
không có vướng víu nối giữa các cửa sổ hay giữa các tầng. Vùng nhìn của mỗi pixel tăng
$3\times3 \to 5\times5 \to 7\times7$. Từ tầng 2, $\phi$ khác 0 nên mã hóa dùng toàn mặt cầu Bloch.
Mỗi tầng $108+56=164$ tham số; ba tầng $492$. Mỗi ảnh cần $256\times3=768$ mạch 9 qubit; mô phỏng statevector
$2^9=512$ biên độ; khi train dùng gradient checkpoint (tính lại trạng thái khi lan truyền ngược).

Đầu ra encoder: $\theta_p=\theta^{(3)}_p$, $\phi_p=\phi^{(3)}_p$, $p=0,\dots,255$; $\theta$ không bị kẹp.

% =====================================================================================================
\section{Mã hóa vào cây và thứ tự quét}

\subsection{256 qubit lá}
Mỗi pixel thành một qubit theo cùng công thức \eqref{eq:enc}:
\begin{equation}
\ket{q_p} = \cos\tfrac{\pi\theta_p}{2}\ket0 + e^{i\pi\phi_p}\sin\tfrac{\pi\theta_p}{2}\ket1,
\qquad \rho_p = \ket{q_p}\bra{q_p}.
\end{equation}
256 qubit ở trạng thái tích. Không có encoder thì $\theta_p=x_p$, $\phi_p=0$: qubit chỉ phụ thuộc một pixel và nằm
trên nửa đường tròn $Y=0$ của mặt cầu Bloch; có encoder thì phụ thuộc vùng $7\times7$ và dùng toàn mặt cầu.

\subsection{Ba thứ tự quét}
Cùng 256 qubit được đưa vào ba cây, mỗi cây sắp lá theo một hoán vị $\pi_t$ của các pixel:
\begin{itemize}[nosep]
\item hàng: $p = 16i+j$ (thứ tự gốc);
\item cột: sắp theo $(j, i)$;
\item Morton (đường chữ Z): sắp theo mã xen bit
$m(i,j) = \sum_{b=0}^{3}\big(j_b\,2^{2b} + i_b\,2^{2b+1}\big)$, với $i_b, j_b$ là bit thứ $b$ của $i$ và $j$.
\end{itemize}
Cây ghép từng cặp lá liên tiếp, nên một qubit ở tầng có $n$ qubit đại diện cho $256/n$ pixel liên tiếp theo thứ tự
quét: với Morton đó là các khối $1\times2$, $2\times2$, $2\times4$, $4\times4$, $4\times8$; với quét hàng là các dải
$1\times2$, $1\times4$, $1\times8$, $1\times16$, $2\times16$.

% =====================================================================================================
\section{Cây QCNN tensor network (mỗi cây $1.608$ tham số)}

\chenhinh{3}{Một cây: 5 tầng conv-pool, đo tầng 32, re-upload, strongly-entangling 8 qubit, đặc trưng Pauli.}

\subsection{Unitary học được}
Mọi unitary của cây được tham số hóa bằng hàm mũ của ma trận Hermite:
\begin{equation}
U = \exp(iH), \qquad H = \tfrac12\big(A + A^\dagger\big),\qquad A \in \CC^{d\times d},
\end{equation}
với $d=2$ (8 số thực) hoặc $d=4$ (32 số thực); phần thực và ảo của $A$ khởi tạo $\mathcal N(0,\,0{,}01^2)$,
nên mọi unitary bắt đầu gần đơn vị.

\subsection{Năm tầng conv-pool: $256\to128\to64\to32\to16\to8$}
Ở tầng $\ell=1,\dots,5$, một cặp unitary 2 qubit $(U_c^{(\ell)},U_p^{(\ell)})$ (64 tham số) dùng chung cho mọi cặp:
\begin{equation}
\rho^{(\ell)}_j = \Tr_2\Big[ V_\ell\,\big(\rho^{(\ell-1)}_{2j}\otimes\rho^{(\ell-1)}_{2j+1}\big)\,V_\ell^\dagger \Big],
\qquad V_\ell = U_p^{(\ell)}U_c^{(\ell)},
\end{equation}
với $\rho^{(0)}_j$ là 256 lá theo thứ tự quét và $\Tr_2$ là lấy vết riêng phần trên qubit thứ hai.
Vì qubit thứ hai bị lấy vết ngay, các qubit trong cùng tầng luôn độc lập: trạng thái của tầng là tích các ma trận
mật độ $2\times2$, và vector Bloch của từng qubit là toàn bộ thông tin của tầng.

\subsection{Đo tầng 32 (\texttt{side-level 32})}
Sau tầng $\ell=3$ còn 32 qubit; lấy vector Bloch của từng qubit:
\begin{equation}
\bm s_t = \big(\Tr[\rho^{(3)}_j X],\ \Tr[\rho^{(3)}_j Y],\ \Tr[\rho^{(3)}_j Z]\big)_{j=0}^{31} \in [-1,1]^{96}.
\end{equation}

\subsection{Re-upload trên 8 qubit cuối}
Qubit cuối $q\in\{0,\dots,7\}$ đại diện cho 32 lá liên tiếp $B_q$. Lấy góc trung bình của các lá đó
(từ độ lớn biên độ) và một hệ số học được $s_q$ (khởi tạo 0, 8 tham số):
\begin{equation}
\bar\vartheta_q = \frac{1}{32}\sum_{p\in B_q} 2\operatorname{atan2}\big(|\sin\tfrac{\pi\theta_p}{2}|,\,|\cos\tfrac{\pi\theta_p}{2}|\big),
\qquad
\rho^{(5)}_q \leftarrow R_Y(s_q\bar\vartheta_q)\,\rho^{(5)}_q\,R_Y(s_q\bar\vartheta_q)^\dagger .
\end{equation}

\subsection{Strongly-entangling trên 8 qubit}
Ghép 8 qubit thành ma trận mật độ đầy đủ $\rho = \bigotimes_{q=0}^{7}\rho^{(5)}_q \in \CC^{256\times256}$.
Bốn lớp $m=0,\dots,3$, mỗi lớp: một unitary 1 qubit $u^{(m)}_q$ trên mỗi qubit (8 tham số mỗi cái), rồi 8 unitary
2 qubit $W^{(m)}_q$ trên cặp $(q,\ (q+r_m)\bmod 8)$ với khoảng cách $r_m = m+1$ (32 tham số mỗi cái):
\begin{equation}
\rho \leftarrow \Big(\prod_{q} W^{(m)}_{q,\,q+r_m}\Big)\Big(\bigotimes_q u^{(m)}_q\Big)\,\rho\,
\Big(\bigotimes_q u^{(m)}_q\Big)^\dagger\Big(\prod_{q} W^{(m)}_{q,\,q+r_m}\Big)^\dagger .
\end{equation}
Mỗi lớp $8\times8+8\times32=320$ tham số; 4 lớp $1.280$. Chỉ ở đây các qubit mới vướng víu với nhau.

\subsection{Đặc trưng Pauli}
\begin{equation}
\bm f_t = \Big( \big(\Tr[\rho\,\sigma^a_q]\big)_{a\in\{X,Y,Z\},\,q}\ ,\ \big(\Tr[\rho\,\sigma^a_q\sigma^b_{q'}]\big)_{a,b\in\{X,Y,Z\},\,q<q'} \Big)
\in [-1,1]^{24+252} = [-1,1]^{276}.
\end{equation}
Mỗi cây cho $[\bm f_t;\bm s_t]\in\RR^{372}$. Tham số mỗi cây: $5\times64 + 8 + 1.280 = 1.608$.

% =====================================================================================================
\section{Head cổ điển ($11.948$ tham số)}
Với mỗi cây $t\in\{\text{hàng},\text{cột},\text{Morton}\}$:
\begin{equation}
\bm h_t = \mathrm{Dropout}\Big(\mathrm{GELU}\big(\mathrm{LN}\big(W_t[\bm f_t;\bm s_t]+\bm b_t\big)\big)\Big) \in \RR^{10},
\qquad W_t\in\RR^{10\times372}\quad(3.750 \text{ tham số kể cả LN}).
\end{equation}
Gộp:
\begin{equation}
\bm z = [\bm h_{\text{hàng}};\bm h_{\text{cột}};\bm h_{\text{Morton}}]\in\RR^{30},\qquad
\text{logit} = W_2\,\mathrm{Dropout}\Big(\mathrm{GELU}\big(\mathrm{LN}(W_1\bm z+\bm b_1)\big)\Big)+\bm b_2 \in\RR^{10},
\end{equation}
$W_1\in\RR^{16\times30}$, $W_2\in\RR^{10\times16}$ ($698$ tham số). Dropout $p=0$ trong các run đã chạy.

% =====================================================================================================
\section{Huấn luyện}
\begin{itemize}[nosep]
\item Hàm mất mát cross-entropy trên logit; Adam, learning rate $10^{-3}$, weight decay $0$, batch 64.
\item ReduceLROnPlateau theo val loss: nhân $0{,}5$ sau 3 epoch không giảm, tối thiểu $10^{-6}$.
\item Tối đa 50 epoch; dừng sớm sau 10 epoch val loss không giảm. Seed 0.
\item Lưu checkpoint val acc cao nhất (MAX) và val loss thấp nhất (CK). Chọn mô hình chỉ bằng validation.
\item Thời gian trên RTX 4090: khoảng $1.350$ s/epoch khi chạy một mình.
\end{itemize}

% =====================================================================================================
\section{Kết quả hiện có (validation, seed 0)}
\begin{table}[htbp]\centering
\caption{MAX val acc (\%). Một seed; chênh dưới $0{,}3$ điểm chưa đủ kết luận. Số có dấu * là giữa run.}
\begin{tabular}{lrr}
\toprule
Mô hình & Tham số & MAX val \\
\midrule
level32, không encoder & 19.118 & 98,84 \\
level32 + encoder CNN deep12p (cửa sổ $7\times7$) & 18.418 & 99,14 \\
level32 + qconv2x2 (cửa sổ $5\times5$) & 16.992 & $\approx$98,8 \\
\textbf{level32 + qconv4x3 (mô hình này)} & \textbf{17.264} & \textbf{99,03}* \\
qconv4x3 + AvgPool + Linear (bỏ cây) & 822 & 95,77* \\
CNN cổ điển cnn5k & 5.578 & 99,11 \\
\bottomrule
\end{tabular}
\end{table}
Còn thiếu: đối chứng encoder CNN cùng cỡ \texttt{deep6p} (17.272), thay cây bằng CNN cổ điển cùng cỡ
\texttt{enc-cnn qconv4x3} (17.238), encoder đóng băng \texttt{qconv4x3f}, thêm seed, test chính thức.

% =====================================================================================================
\section{Ghi chú}
\begin{itemize}[nosep]
\item Mọi kỳ vọng được tính chính xác (không lấy mẫu, không nhiễu). Encoder là các mạch 9 qubit độc lập; trong cây
các qubit độc lập giữa các tầng, chỉ 8 qubit cuối vướng víu. Do đó mô hình mô phỏng cổ điển hiệu quả.
\item Mã hóa góc một qubit mỗi pixel theo $\cos\frac{\pi x}{2}, \sin\frac{\pi x}{2}$ đã dùng trong
Stoudenmire \& Schwab (2016); mã hóa hai góc mỗi qubit tương tự ``dense angle encoding'' (LaRose \& Coyle, 2020);
mạch trượt trên ô ảnh rồi đo tương tự quanvolution (Henderson và cộng sự, 2019). Cần kiểm lại trích dẫn trước khi
dùng.
\end{itemize}

\end{document}
